% Generated draft: requires mathematical and source review.
\documentclass[11pt]{article}
\usepackage[margin=25mm]{geometry}
\usepackage{fontspec}
\setmainfont{Latin Modern Roman}
\usepackage{amsmath,amssymb,mathtools,unicode-math}
\setmathfont{Latin Modern Math}
\usepackage{enumitem,needspace,xurl,hyperref}
\setlength{\parindent}{0pt}
\setlength{\parskip}{5pt}
\setlength{\emergencystretch}{3em}
\newcommand{\N}{\mathbb{N}}
\newcommand{\Z}{\mathbb{Z}}
\newcommand{\Q}{\mathbb{Q}}
\newcommand{\R}{\mathbb{R}}
\newcommand{\C}{\mathbb{C}}
\newcommand{\F}{\mathbb{F}}
\newcommand{\A}{\mathbb{A}}
\newcommand{\PP}{\mathbb{P}}
\newcommand{\E}{\mathbb{E}}
\newcommand{\eps}{\varepsilon}
\newcommand{\dd}{\,\mathrm{d}}
\begin{document}
\section*{UnsolvedMath ID: 30002916}
\textbf{Source problem code:} catalogue entry
\subsection{Definitions and mathematical statement}
The source asks a qualitative question and does not supply enough hypotheses to determine a unique initial-value problem. The following definitions describe the conventional setting, with the missing qualifications made explicit.

Let \(X\) be a complex manifold. In the usual interpretation of Kähler--Ricci flow from currents, the initial datum \(T_0\) is a closed positive current of bidegree \((1,1)\). Positivity means that \(T_0\) evaluates nonnegatively on compactly supported smooth strongly positive test forms of complementary bidegree. Closedness and bidegree are conventional here; the supplied statement says only ``positive current.''

Locally write \(T_0=(i/\pi)\partial\bar\partial u\), where \(u\) is a plurisubharmonic potential. Using this normalization, its Lelong number at \(x\in X\) is
\[
\nu(T_0,x)=\liminf_{z\to x}\frac{u(z)}{\log\lVert z-x\rVert},
\]
computed in local holomorphic coordinates centered at \(x\). This number is independent of the coordinates and the choice of local potential. Write
\[
E_+(T_0)=\{x\in X:\nu(T_0,x)>0\}.
\]
The phrase ``has positive Lelong numbers'' is interpreted here as \(E_+(T_0)\ne\varnothing\); the source gives no more specific condition on this locus.

For a smooth family of Kähler forms \(\omega(t)\), the unnormalized Kähler--Ricci flow equation is
\[
\frac{\partial\omega(t)}{\partial t}=-\operatorname{Ric}(\omega(t)),
\qquad
\operatorname{Ric}(\omega)=-i\partial\bar\partial\log\det(g_{j\bar k}),
\]
where \((g_{j\bar k})\) is the associated Hermitian metric matrix. A current-valued evolution starting from \(T_0\) would ordinarily have the initial trace
\[
T(t)\rightharpoonup T_0\quad\text{as }t\downarrow0,
\]
meaning convergence against every compactly supported smooth test form of complementary degree. Neither this initial-trace convention nor a definition of weak solution is supplied in the record; the smooth equation alone does not define the Ricci form of an arbitrary singular current.

\textbf{Source question.} What is the behavior of Kähler--Ricci flow with initial positive current \(T_0\) having positive Lelong numbers? In the conventional interpretation above, the question concerns weak evolutions with \(E_+(T_0)\ne\varnothing\). The source does not specify the ambient manifold, compactness, initial cohomology class, normalized versus unnormalized flow, time interval, weak-solution concept, or a particular behavior to prove. Thus it supplies a research question rather than a uniquely quantified conjecture; the displayed conventions do not resolve those omissions.
\subsection{Short English statement}
The problem asks how Kähler--Ricci flow behaves when its initial current has positive Lelong numbers, which measure local singularity concentration. The source leaves the geometric hypotheses, weak-solution framework, and intended conclusion unspecified, so a unique precise conjecture cannot be extracted without further source clarification.
\subsection{Sources}
\begin{itemize}
\item Local UnsolvedMath registry: permanent ID \(30002916\), exact source problem code \texttt{catalogue entry}, title ``Kähler--Ricci Flow from Currents with Positive Lelong Numbers,'' in the Oberwolfach Reports Open Problems collection. No individual proposer is supplied; the registry gives proposed year \(2015\).
\item Supplied report citation: \emph{Differentialgeometrie im Großen (2016)}, Oberwolfach Reports, DOI \texttt{10.4171/owr/2015/31}, pp.\ 1759--1807. \url{https://doi.org/10.4171/owr/2015/31}. The record describes these reports as documenting mathematical workshops and their open-problem discussions. The title year \(2016\), DOI year component \(2015\), and proposed year \(2015\) are retained as supplied; their relationship is unverified.
\item Supplied literature reference, labeled ``Di Nezza and Lu primary paper'': \url{https://doi.org/10.1016/j.aim.2016.10.011}. No paper title or specific result is supplied.
\item Supplied literature reference, labeled ``Di Nezza Guedj and Lu primary paper'': \url{https://arxiv.org/abs/2606.15816}. No paper title or specific result is supplied.
\item Literature assessment dated exactly \texttt{2026-08-21}: the supplied curation labels the problem ``partially solved'' and reports that weak Kähler--Ricci flow from positive currents is understood in several regularity and uniqueness settings, including divisorial singularities, while the full positive-Lelong-number behavior remains unsettled. This is supplied dated curation derived from the OWR record, not an independently verified assessment. The record explicitly requests independent verification of time-sensitive status and cited claims.
\item The registry labels the statement \texttt{corrected\_verified} and says that mathematical content was preserved while notation, typography, or omitted definitions were repaired for self-contained publication. These are supplied editorial assertions, not independent mathematical verification of this draft.
\end{itemize}
\end{document}
